\section{Decision one: which liturgical day is it?}

A liturgical day is a day set apart for the liturgy, running midnight to
midnight, whose celebration ordinarily runs from Matins to Compline (\rg{4},
\rg{5}). On each day the Office and Mass are ``aut de dominica, aut de feria,
aut de vigilia, aut de festo, aut de octava'' (\rg{6}). Those five are the
whole inventory; there is nothing else a day can be, except that a sixth
possibility, the Office of Our Lady on Saturday, occupies fourth-class
Saturdays (\rg{78}).

So the first task is mechanical: write down every one of those five things
that falls on your date, in every calendar that binds you.

\subsection{Which calendar binds you}

\rgm{274} settles this before anything else: ``Missa dicenda est iuxta
calendarium aut ecclesiae vel oratorii in quo Missa celebratur, aut loci, aut
ipsius sacerdotis celebrantis, aut Ecclesiae universae, prout infra
exponitur.'' The cases that follow are not interchangeable.

\begin{tabletwo}{Where Mass is said}{Calendar to be followed}
In a church or public oratory & That church's or oratory's own calendar, for
every priest, diocesan or religious alike; likewise in the principal
semi-public oratory of a seminary, religious house, college, hospital or
prison (\rgm{275}).\\
In a secondary oratory of such a house & Either that oratory's calendar or the
priest's own (\rgm{276}).\\
In a private oratory, or on a portable altar outside a sacred place & Either
the calendar of the place (\rg{53}\,a) or the priest's own (\rgm{277}).\\
In an oratory fixed in a ship & The calendar of the universal Church; a priest
celebrating outside it on a portable altar, or lawfully in the course of air,
river or rail travel, may use either the universal calendar or his own
(\rgm{279}).\\
In a church of a different rite & That church's calendar for feasts, their
grade, commemorations and the collect imposed by the Ordinary; but the
variable parts are taken from that rite, with its own Ordinary and ceremonies
(\rgm{284}).\\
\end{tabletwo}

\rgm{278} then overrides the freedom: whatever calendar a priest might
otherwise follow, he \emph{must} say the Mass of the principal Patron of the
nation, region or province, diocese, town or city, of the anniversary of the
dedication of the cathedral, and of other days actually observed as feasts of
precept. Which particular feasts these are is a question for the approved
proper, not for this book.

\subsection{The building blocks, with their classes}

\subsubsection*{Sundays}

Sundays are of the first or second class (\rg{10}). First-class Sundays are
exactly six kinds: the four of Advent, the four of Lent, the two of
Passiontide, Easter, Low Sunday and Pentecost (\rg{11}); Easter and Pentecost
are also first-class feasts with octaves. \emph{Every other Sunday is second
class} (\rg{12}). The Sunday Office begins at First Vespers on the preceding
Saturday (\rg{13}).

Two rules about Sundays do a great deal of work later:

\begin{rulebox}{\rg{14}}
Dominica celebratur suo die, iuxta rubricas. Officium et Missa dominicae
impeditae nec anticipantur nec resumuntur.
\end{rulebox}

\begin{rulebox}{\rg{17} (opening)}
Dominica excludit, per se, assignationem perpetuam festorum.
\end{rulebox}

\noindent
An impeded Sunday is simply gone that year --- it is not moved forward, back,
or into the following week. And no feast may be perpetually assigned to a
Sunday except the five the rubric itself names: the Holy Name, the Holy
Family, Trinity Sunday, Christ the King, and those first-class feasts of the
Lord already assigned to a second-class Sunday in a particular calendar. Those
five ``locum tenent dominicae occurrentis cum omnibus iuribus et privilegiis;
de dominica, proinde, nulla fit commemoratio'' (\rg{17}).

\subsubsection*{Ferias}

A feria is any weekday other than Sunday (\rg{21}), and ferias run through all
four classes. This is where most collisions with saints' days are actually
decided, and where the code most sharply departs from the habits of an older
missal.

\begin{tablerule}{Feria}{What it does}{Locus}
First class: Ash Wednesday; all the ferias of Holy Week & Take precedence over
any feast whatever, and admit no commemoration except one privileged one. &
\rg{23}\\
Second class: the ferias of Advent from 17 to 23 December; the Ember Days of
Advent, of Lent and of September & Take precedence over \emph{particular}
second-class feasts; if impeded, must be commemorated. & \rg{24}\\
Third class (a): the ferias of Lent and Passiontide from the Thursday after
Ash Wednesday to the Saturday before Passion Sunday inclusive, not named above
& \textbf{Take precedence over third-class feasts.} If impeded, must be
commemorated. & \rg{25}\,a\\
Third class (b): the ferias of Advent to 16 December inclusive, not named
above & \textbf{Yield to third-class feasts.} If impeded, must be
commemorated. & \rg{25}\,b\\
Fourth class: every feria not named above & Never commemorated at all. &
\rg{26}\\
\end{tablerule}

The two third-class rows are printed as adjacent clauses of a single rubric
and they say opposite things. A third-class saint on a Lenten weekday loses
and is commemorated; the same saint on an Advent weekday before 17 December
wins and the feria is commemorated. Nothing in the class numbers predicts
this; only \rg{25} does.

\subsubsection*{Vigils}

A vigil is a liturgical day preceding a feast and preparing for it (\rg{28});
the Easter Vigil is excepted, ``cum non sit dies liturgicus,'' and is
celebrated in its own manner. First-class vigils are two: Christmas, which
takes the place of the Fourth Sunday of Advent with no commemoration of it,
and Pentecost; they yield to nothing and admit no commemoration (\rg{30}).
Second-class vigils are four: the Ascension, the Assumption, the Nativity of
St John the Baptist, and Ss.\ Peter and Paul; they outrank third- and
fourth-class days and are commemorated if impeded (\rg{31}). The one
third-class vigil is that of St Lawrence (\rg{32}). And a second- or
third-class vigil is dropped altogether --- ``penitus omittitur'' --- if it
falls on any Sunday or on a first-class feast, or if the feast it precedes is
itself transferred or reduced to a commemoration (\rg{33}).

\subsubsection*{Octaves}

\begin{rulebox}{\rg{64}}
Celebrantur tantum octavae Nativitatis Domini, Paschatis et Pentecostes,
exclusis omnibus aliis, sive in calendario universali sive in calendariis
particularibus.
\end{rulebox}

\noindent
Three, and no more, anywhere. The octaves of Easter and Pentecost are first
class and their days are first-class days; the octave of Christmas is second
class, its days second class, its octave day first class (\rg{66}--\rg{67}).
The Christmas octave is then arranged by name, day by day, in \rg{68}: St
Stephen, St John, the Holy Innocents (all second-class feasts), a
commemoration of St Thomas of Canterbury on 29 December, a commemoration of St
Sylvester on 31 December, and from particular calendars only first-class
feasts of saints inscribed in the universal calendar on those same days, the
rest being transferred past the octave.

\subsubsection*{Feasts, and where a feast comes from}

Feasts are of the first, second or third class (\rg{36}), and are either
universal --- inscribed by the Holy See in the calendar of the universal
Church, and to be celebrated by everyone of the Roman rite (\rg{39}) --- or
particular, that is proper or indulted (\rg{38}). The proper feasts a
particular calendar must carry are enumerated exhaustively in
\rg{42}--\rg{46}, and their classes are fixed there, not chosen locally:

\begin{tablerule}{Proper feast of}{Feast, with the class the rubric fixes}{Locus}
A nation, region or province & Principal Patron (I class); secondary Patron
(II class). & \rg{42}\\
A diocese or equivalent territory & Principal Patron (I); anniversary of the
dedication of the cathedral (I); secondary Patron (II); saints and blessed
duly inscribed in the Martyrology having a special relation of origin, long
residence or death (II or III class, or a commemoration). & \rg{43}\\
A place, town or city & Principal Patron (I); secondary Patron (II). &
\rg{44}\\
A church or public or semi-public oratory taking a church's place & Anniversary
of dedication, if consecrated (I); feast of the Title, if consecrated or at
least solemnly blessed (I); the saint whose body is kept there (II); a blessed
whose body is kept there (III). & \rg{45}\\
An Order or Congregation & Title (I); canonised Founder (I) or beatified
Founder (II); principal Patron of the whole Order, and of each religious
province (I); secondary Patron (II); saints and blessed who were its members
(II or III, or a commemoration). & \rg{46}\\
\end{tablerule}

Only first-class feasts have a right of transfer when accidentally impeded
(\rg{95}); everything else in this table is either commemorated or dropped for
the year. That is worth noticing now, because it means a diocesan
second-class Patron is materially weaker than the number ``II'' suggests.

\subsubsection*{Our Lady on Saturday}

On Saturdays carrying a fourth-class feria, the Office is of Our Lady on
Saturday (\rg{78}). It is the twenty-seventh of the twenty-eight rows of the
precedence table --- a fourth-class \emph{Office}, not a feast --- and it
begins at Matins and ends after None (\rg{79}). Its Mass therefore exists and
may be sung with the solemn tone (\rgm{515}\,b), but it is displaced by any
occurring feast whatever.

\subsection{The seasons, and why they matter here}

The code fixes the seasons by their exact boundaries (\rg{71}--\rg{77}):
Advent from First Vespers of Advent Sunday to None of the vigil of Christmas;
Christmastide to 13 January, subdivided at Epiphany; Septuagesima from First
Vespers of Septuagesima Sunday to Compline of Shrove Tuesday; Lent from Matins
of Ash Wednesday to the Mass of the Easter Vigil exclusive, subdivided at
Passion Sunday; Paschaltide from the beginning of the Easter Vigil Mass to
None of the Saturday in the octave of Pentecost, subdivided at the Ascension
and the vigil of Pentecost; and \latin{tempus ``per annum''} from 14 January
to the Saturday before Septuagesima and from First Vespers of Trinity Sunday
to None of the Saturday before Advent.

These are not decorative. Five of the assembly rules in Section~6 and almost
all the Preface rules in Section~7 are keyed to a season boundary rather than
to the day's class or the Mass's formulary, and a boundary is either crossed
or it is not.
